% !TEX TS-program = lualatex
%\documentclass[handout]{beamer}
\documentclass{beamer}
\input{../common/misomva}
\begin{document}
% Topic-specific title slide

\newcommand{\topictitle}{Google PageRank: Computation}

\newcommand{\topicsubtitle}{Power Method and Large-Scale Implementation}
\newcommand{\misoclasstinytitleslide}{Partially based on material by:  Andreas Krause, \\[-1em] Branislav Kveton, Michael Kearns}

\MakeTopicSlide

\begin{frame}

\frametitle{Success story \#2 \emphcol{Google \pagerank}}

\textbf{History:} {\small [Desikan, 2006]} 
\begin{itemize}
\item The anatomy of a large-scale hypertextual web
search engine [Brin \& Page 1998] \pause
\item  US patent for PageRank granted in 2001\pause
\item  Google indexes 10's of billions of web pages
(1 billion = $10^9$) \pause
\item  Google serves $\geq 200$ million queries per day \pause
\item  Each query processed by $\geq 1000$ machines \pause
\item  All search engines combined process more
than 500 million queries per day
\end{itemize}


\end{frame}
\begin{frame}

\frametitle{Success story \#2 \emphcol{Google \pagerank}}

\begin{block}{}
\emph{Problem:} Find an eigenvector of a stochastic matrix.  
\end{block}
 \pause
\begin{itemize}
\item  $n = 10^9$ !!! \pause
\item luckily: \textbf{ sparse} \pause (average outdegree: 7) \pause
%\item better than a simple centrality measure (e.g., degree) \pause
\item power method \pause
\end{itemize}
\vspace{-1.5em}
\begin{align*} % requires amsmath; align* for no eq. number
   \onslide<6->{\bv_0 &= \left(1_A\quad 0_B\quad  0_C\quad  0_D\quad  0_E\quad  0_F\right)\transpose \\}
   \onslide<7->{\bv_1 &=  \bG  \bv_0 \\} 
   \onslide<8->{\bv_{t+1} &= \bG \bv_{t} } % \onslide<9->{= \bG^{t+1}\bv} 
\end{align*}
\onslide<9->{$\bv_{t+1} = \bv_{t} \implies \bG\bv_{t} = \bv_{t}$}
\onslide<10->{\notecol{ \quad \tiny  and we found the steady vector}}
\onslide<11->{\rightquestionbox{But wait,  $\bM$ is sparse, but $\bG$ is dense! What to do?}} 
\onslide<12-|handout:0>{
\begin{flushright}
\vspace{-1em}
\notecol{ \tiny we store only $\bM$ but do computations as with $\bG$ }
\end{flushright}
}
\end{frame}
%%%%%%%%%%%%%%%%%%%%%%%%%%%%%%%%%%%%%%%%%%%%%%%%%%%%%%%%%


\MakeLastSlide
\end{document}
